\chapter{Discussion}
\label{ch:discussion}

% ---------------------------------------------------------------------------
% CHAPTER 4 WRITING OUTLINE (draft prose in a later iteration)
% ---------------------------------------------------------------------------
%
% sec:discussion-analysis — Analysis of the Results
%   • MACS-18 vs ASOCA: CESR improvement (95.62% vs 83.19%) validates H1;
%     richer topology explains higher MPE and lower OISR (H2).
%   • Runtime increase on MACS-18 (9.60 min vs 5.78 min) confirms H3;
%     still within good band vs manual 30–120 min review.
%   • ESR 100% on both cohorts: pipeline operational reliability.
%   • Synthetic validation: max %AS error 1.31% (<5% inter-observer band);
%     confirms Block 2 logic is sound when ground truth is known.
%   • SUS mean 70.83 (Good band): pilot usability acceptance; n=3 limits
%     generalizability; Q9 neutral/confident split supports decision-support role.
%
% sec:discussion-limitations — Technical and Anatomical Limitations
%   • CRITICAL: systematic stenosis overestimation on real patients
%     (fixed W=10 mm, bifurcations, partial-volume, segmentation artifacts).
%   • No clinical stenosis/CAD-RADS ground truth evaluated — by design.
%   • Ostium detection degrades on MACS-18 despite better centerlines.
%   • Manual visual audit for OISR/CESR/EISR — subjective, not fully blinded.
%   • Single-center, small SUS sample; no longitudinal clinical outcome data.
%
% sec:discussion-future — Future Work (from cardiologist feedback, Results v2 §4.2)
%   4.2.1 Algorithmic and metric refinements:
%     • Dynamic clinician-adjustable proximal/distal reference windows.
%     • Adaptive centerline resolution (~2 mm default; finer in diseased zones).
%     • Plaque quantification alongside stenosis (clinical priority).
%   4.2.2 Visualization and interface enhancements:
%     • Continuous lumen cross-sectional visualization along centerline.
%     • CMPR and raw CCTA overlay for ground-truth validation in dashboard.
%     • Aggregated affected-zone highlighting (e.g. regions ≥50% stenosis).
%     • Interactive threshold adjustment and explainability features.
%   4.2.3 Clinical workflow integration:
%     • Deeper integration with patient prioritization system (ferrer2024).
%     • Pop-up rapid-assessment within prioritized patient lists at HSP.
%     • Scale SUS study to larger cardiologist cohort; structured UX follow-up.
% ---------------------------------------------------------------------------

\section{Analysis of the Results}
\label{sec:discussion-analysis}

% Content to be written in a later iteration.

\section{Technical and Anatomical Limitations}
\label{sec:discussion-limitations}

% Content to be written in a later iteration.

\section{Future Work}
\label{sec:discussion-future}

% Content to be written in a later iteration.
